\section{Hạn chế của khoa học}

\subsection{So sánh khoa học với nghệ thuật}

\

Hiển nhiên rằng, không phải chỉ nguyên khoa học là bộ môn đưa việc phản ánh thế giới làm trọng tâm. Thiên nhiên và con người liên tục được phản chiếu trong các bài văn thơ, bài hát hay những bức tranh. Hơn thế nữa, trong những tác phẩm nghệ thuật đó, ngoài việc chúng ta được cung cấp không gian và thời gian của sự vật, chúng ta còn nhìn thấy suy nghĩ nội tâm của các nhân vật, đặt mình vào trong dòng đời của người kể hoặc người được kể, để từ đó, chúng ta có thể trải nghiệm những khía cạnh của cảm xúc mà có thể chúng ta chưa từng được biết. Trái lại, tác phẩm khoa học thì khó có thể bộc bạch tinh thần của con người như vậy. Tác giả thường được nghe rằng cần phải trình bày các nghiên cứu như một ``câu chuyện khoa học'', nhưng tác giả cảm thấy câu chuyện này tương đối nhạt. Về mặt kết cấu, các tác phẩm theo văn phong khoa học vẫn có một kết cấu giống như các bài tự sự: có mở đầu, có bối cảnh, có hành động, có kết quả. Tuy nhiên, về mặt nội dung, việc thể hiện cảm xúc chủ quan là một điều cấm kị, do biểu cảm là thứ không cân đo đong đếm được. Nhiệm vụ của một người viết ra khám phá là phải làm sao cho người đọc, với một lượng kiến thức chuyên môn sâu nhưng không đến trình độ chuyên gia, có thể nhanh chóng hiểu và kiểm tra được tính chính xác của bài viết. Điều này dẫn đến sự thiếu hụt của tính ``người'', biểu hiện qua sự hạn chế tới mức tối đa của các ngôn ngữ biểu cảm và các biện pháp tu từ trong các công trình khoa học được viết ra. Thực vậy, các sách sử học có thể chứng minh ảnh hưởng tiêu cực lên người dân Việt Nam dưới thời kì chịu đựng áp bức của thực dân Pháp, nhưng liệu các số liệu và sự kiện trong những cuốn sách đó có tác động mạnh mẽ đến trái tim người đọc như các sáng tác hiện thực phê phán như ``Tắt đèn'' của Ngô Tất Tố, ``Bước đường cùng'' của Nguyễn Công Hoan, và ``Chí Phèo'' của Nam Cao được hay không? Nếu bạn đọc chưa từng tiếp cận những tác phẩm này, thì hãy tìm đọc chúng, và tự đưa ra câu trả lời cho mình.

\subsection{So sánh khoa học với chủ nghĩa duy tâm}

\

Tuy khoa học và văn học là hai phạm trù tương đối khác nhau, chúng vẫn ít nhiều sử dụng thế giới vật chất để quyết định ra điều kiện và hoàn cảnh mà con người có thể hoạt động. Thế còn ở một góc nhìn khác thì sao, khi quan niệm rằng ý thức (có thể không phải của con người) không chỉ phản ánh mà còn có khả năng xây dựng các sự vật và hiện tượng tự nhiên và xã hội? Thì khi đó, chào mừng bạn đọc đến với tư tưởng chủ nghĩa duy tâm, và ở một cấp độ cao nhất, bạn đọc sẽ thấy ở đó những tôn giáo cùng nhau hoạt động. Cần nhìn nhận rằng, mặc dù khoa học và tôn giáo nằm ở hai thái cực đối lập nhau, nhưng không có nghĩa là chúng ta chỉ chấp nhận một phương diện và phủ nhận phương diện còn lại. Ánh sáng, mà sau này chúng ta sẽ tìm hiểu, có lưỡng tính sóng hạt. Giải thích rõ ra, chúng ta có thể coi ánh sáng như là những hạt quang tử (hay hạt phô-ton) ở trong một số trường hợp và như là sóng ánh sáng ở những trường hợp khác. Với người ít hiểu biết về vật lí, sóng và hạt có vẻ như là hai khái niệm tách biệt khỏi nhau, nhưng với những người có kiến thức sâu hơn, họ sẽ thấy rằng cả hai khái niệm của ánh sáng giúp cho họ cái nhìn đầy đủ hơn. Giống như vậy, một người ít hiểu biết hoặc với những tư tưởng cực đoan sẽ cho rằng khoa học và đức tin không thể ngồi chung một chỗ. Nhưng với chúng ta, những người có góc nhìn toàn diện hơn, hoàn toàn có thể nhận thấy rằng khoa học không có khả năng đáp ứng \emph{mọi} nhu cầu của mỗi cá nhân, và từ đó chấp nhận tính nhị nguyên của xã hội loài người.